\documentclass[10pt,a4paper]{amsart}
\usepackage[margin=19mm]{geometry}
\usepackage{amsmath,amssymb,mathtools}
\usepackage[scr=rsfs,cal=euler]{mathalfa}
\usepackage{xfrac}
\usepackage[unicode,hidelinks,bookmarks=false]{hyperref}
\newcommand{\ie}{i.e.\ }
\newcommand{\cf}{cf.\ }
\newcommand{\resp}{respectively\ }
\newcommand{\Subsection}{\subsection}
\providecommand{\nobreakdash}{\nobreak\hbox{-}}
\setlength{\emergencystretch}{2em}
\multlinegap0pt
\def\v{{\bf v}}
\def\R{{\cal R}}

\newtheorem{definition}{Definition}
\newtheorem{proposition}{Proposition}
\newtheorem{corollary}{Corollary}
\newtheorem{theorem}{Theorem}
\newtheorem{lemma}{Lemma}
\newtheorem{remark}{Remark}
\newtheorem{question}{Question}
\newtheorem{fact}{Fact}
\title[On perfect matchings, edge-colourings and eigenvalues of cubic graphs]{On perfect matchings, edge-colourings and eigenvalues of cubic graphs: spectral characterisation of truncated cubic graphs (Proposition 3)}
\author{Willem H. Haemers}
\date{Source: arXiv:2601.03778v3 (CC BY 4.0)}
\begin{document}
\maketitle

\noindent\emph{Source and licence.} On perfect matchings, edge-colourings and eigenvalues of cubic graphs. Version arXiv:2601.03778v3 (CC BY 4.0), distributed under the Creative Commons Attribution 4.0 International licence, http://creativecommons.org/licenses/by/4.0/, as recorded in the arXiv licence metadata for this exact version and shown on the abstract page https://arxiv.org/abs/2601.03778v3. Full licence notice: licenses/arxiv-2601.03778-CC-BY-4.0.txt.\par\smallskip

\noindent\textit{Editorial scope.} Counted unit: the article's Proposition 3 (source0.tex lines 177--179, no source label), ``Being a truncated cubic graph is a property characterised by the spectrum'', delivered with the article's complete original proof (source0.tex lines 181--192, one \texttt{proof} environment of 12 lines, adjacent to the statement, with no gap and no deferral). No mathematics is written, repaired or re-derived: statement and proof are the article's own text line for line, in the article's own notation ($G$, $H$, $T(G)$, $\lambda_i$, $n$), and no line of the counted statement or of the counted proof is omitted, substituted or re-flowed. Required context retained uncounted: the article's own definition of truncation (lines 89--90); the article's own truncation diagram, one figure environment drawn with the standard LaTeX picture environment (lines 92--113, source label \texttt{truncation}), which the definition sentence refers to and whose picture code is therefore retained and delivered inline; the article's definition of the full truncation $T(G)$ and its identification with the line graph of the complete subdivision (lines 115--118); and the article's own Theorem 1 on the spectrum of $T(G)$ (lines 120--127, source label \texttt{trunc-eig}), the statement the counted proposition and its proof refer to as ``the spectrum of $T(G)$ given above''; that theorem carries the article's citation \texttt{ZCC}, which is transcribed character for character in the wrapper's bibliography (source0.tex lines 343--346). Editorial formatting only, in addition: an AMS \texttt{amsart} preamble that restates the article's own macro definitions (\texttt{\textbackslash v}, \texttt{\textbackslash R}) and the theorem environments the retained lines use; sequential numbering of the retained environments, so that the single retained theorem and the single retained proposition are numbered 1 and 1 in their own sequences, whereas the article prints them as Theorem 1 and Proposition 3; the article's own page-geometry settings are replaced by the wrapper's AMS geometry. Exactly one bibliography entry is delivered, the cited key \texttt{ZCC}; no other reference is cited inside a retained span. No graphics-inclusion command occurs inside any retained span, so no external image file is needed or shipped: the truncation diagram is the article's own picture-environment code, and the article's second figure, which uses \texttt{\textbackslash epsfig} to include a JPEG file that the e-print does not contain, lies outside every retained span and is not delivered.

\section*{Retained context: the truncation operation}

Suppose $G$ is a cubic graph. 
{\em Truncation} is the operation that replaces a vertex $v$ of $G$ by a triangle and connects each neighbour of $v$ to one vertex of the triangle (see Fig~\ref{truncation}). 


\begin{figure}[ht]
%\begin{center}
\setlength{\unitlength}{3pt}
\begin{picture}(0,20)(-52,-10)
\thicklines
\put(0,0){\line(-4,3){8}}
\put(0,0){\line(4,3){8}}
\put(0,-8){\line(0,1){8}}
\put(27,-8){\line(0,1){4}}
\put(27,-4){\line(3,5){5}}
\put(27,-4){\line(-3,5){5}}
\put(22,4){\line(1,0){9}}
\put(22,4){\line(-2,1){4}}
\put(32,4){\line(2,1){4}}
\put(10,-2){\huge $\rightarrow$}
\put(0,0){\circle*{2}}
\put(27,-4){\circle*{2}}
\put(22,4){\circle*{2}}
\put(32,4){\circle*{2}}
\end{picture}\caption{Truncation}\label{truncation}
%\end{center}
\end{figure}


If we truncate with respect to every vertex we obtain the truncation of $G$ denoted by $T(G)$.
The truncated graph $T(G)$ can also be described as the line graph of the complete subdivision of $G$.
Using this the spectrum of $T(G)$ can be expressed in terms of the spectrum of $G$;
see Theorem 2.1 in Zhang, Chen and Chen~\cite{ZCC}.


\section*{Retained context: the spectrum of $T(G)$}

\begin{theorem}\label{trunc-eig}\cite{ZCC}
Let $G$ be a cubic graph of order $n$ with spectrum $\lambda_1=3\geq\ldots\geq\lambda_n$.
Then the eigenvalues of the truncated graph $T(G)$ are 
$\frac{1}{2}\pm\frac{1}{2}\sqrt{13+4\lambda_i}$ for 
$i\in\{1,\ldots,n\}$,
%$i=1,\ldots,n$, 
$-2$ and $0$, both with multiplicity $n/2$.
\end{theorem}


\section{The counted unit: Proposition 3 and its complete original proof}

\begin{proposition}
Being a truncated cubic graph is a property characterised by the spectrum.
\end{proposition}

\begin{proof}
Suppose $H$ is a graph with the spectrum of $T(G)$ given above.
Then $H$ is a cubic graph of order $3n$ with $n$ triangles.
Moreover, $H$ has the same number of closed 4-walks as $T(G)$.
But $T(G)$ has no 4-cycles and therefore only trivial closed 4-walks.
So also $H$ has no 4-cycles, which implies that all triangles are 
disjoint, and that two triangles are connected by at most one edge.
Now we define the graph $G'$ whose vertices are the triangles, 
where triangles are adjacent whenever they are connected by an edge in 
$H$. 
Then clearly $H=T(G')$.
\end{proof}

\begin{thebibliography}{99}
\bibitem[ZCC]{ZCC} Fuji Zhang, Yi-Chiuan Chen and Zhibo Chen, {\em Clique-inserted-graphs and spectral dynamics of clique-inserting}, Journal of Mathematical Analysis and Applications 349 (2009), 211--225.
\end{thebibliography}
\end{document}
